% Sumber angka: results/eval_prediksi/summary_all_horizons.csv (17 Agustus 2026).
% Ketiga model dilatih pada SATU jalan yang sama lewat scripts/eval_rf_lstm.py, sehingga
% pembagian data, penskalaan, dan himpunan uji identik bagi ketiganya.
%
% FORMAT (25 Agustus 2026): tabular -> tabularx selebar \textwidth. Sebelumnya keluar
% margin 55,5 pt karena kolom "Clarke A+B (\%)" pada kepala tabel tidak pernah membungkus.
% Satuan dipindah ke keterangan tabel supaya kepala kolom memendek.
\begin{table}[htbp]
\centering
\caption{Kinerja prediksi glukosa pada data uji \emph{hold-out} untuk horizon $+30$ dan $+60$ menit. \emph{Gradient boosting} adalah model produksi; \emph{random forest} dan LSTM adalah pembanding. RMSE dan MAE dalam mg/dL, sedangkan MAPE, Clarke~A, dan Clarke~A+B dalam persen.}
\label{tab:kinerja-prediksi}
\small
\begin{tabularx}{\textwidth}{@{}lLrrrrr@{}}
\toprule
\textbf{Horizon} & \textbf{Model} & \textbf{RMSE} & \textbf{MAE} & \textbf{MAPE} & \textbf{Clarke A} & \textbf{Clarke A+B} \\
\midrule
$+30$ menit & \textbf{Gradient boosting} & 22,25 & 14,72 & 11,13 & 85,35 & 94,70 \\
$+30$ menit & Random forest & 22,60 & 15,10 & 11,53 & 84,63 & 94,35 \\
$+30$ menit & LSTM & \textbf{21,98} & \textbf{14,55} & \textbf{10,93} & \textbf{85,93} & \textbf{94,72} \\
\midrule
$+60$ menit & \textbf{Gradient boosting} & \textbf{33,82} & \textbf{24,41} & \textbf{18,87} & \textbf{68,42} & \textbf{87,32} \\
$+60$ menit & Random forest & 34,24 & 24,83 & 19,38 & 67,89 & 86,94 \\
$+60$ menit & LSTM & 34,14 & 24,62 & 18,87 & 68,40 & 87,22 \\
\bottomrule
\end{tabularx}

\vspace{0.4em}
{\tabnotestyle Cetak tebal menandai nilai terbaik pada tiap horizon. Model produksi \textbf{tidak} unggul pada seluruh sumbu: pada horizon $+30$ menit LSTM sedikit lebih baik, dengan selisih RMSE 0,27~mg/dL dan selisih zona aman 0,02 poin persen, keduanya jauh di bawah batas kebermaknaan klinis. Dasar pemilihan model produksi diuraikan pada Subbab~\ref{sec:eval-prediksi} dan tidak bertumpu pada galat prediksi.\par}
\end{table}
